% Nouveau EGA V (cote 25): archive cote 25, site page 93 (PDF page 94).
% typescript. Draft compared: none.
% AI-transcribed by Claude Opus 5.5 (claude-opus-5-5, adaptive thinking, effort high), Message Batches API, batch ega25b (msgbatch_01A7n5YuXubUPNK1zsepi8RG), 2026-11-13; checked against both supplied scan parts; see CHECK_DECISIONS.tsv.
\scanpage{25/93}
\mspage{87}
de G, on peut également prendre un point à valeurs dans un \var{\inserted[typed]{S-}}{S-}{insertion}{ega25_25_093-v0001}préschéma T,
on obtient alors $\underline{P}_{\hand{\unsure{\xi|\zeta}}} \hand{\longrightarrow} \unsure{P|B}_T$ (T-morphisme, qui est une immersion ferm\unsure{ée|é}
dans le cas envisagé d'abord).

Si $f:X \hand{\longrightarrow} P$ est un morphisme, on pose $\underline{X} = Xx_P\underline{P}$ , et on obtient
un diagramme de même type que le carré précédent
\diagram{\[\begin{array}{ccc}
X & \hand{\unsure{\longleftarrow|\relbar}} & \underline{X} \\
\hand{\downarrow} & & \hand{\downarrow} \\
S & \hand{\longleftarrow} & G
\end{array}\]}
.

Il est alors évident que toutes les questions étudiées dans les n$^{\text{os}}$ 2 à 8
gardent un sens dans le contexte général qu'on vient de poser, et il y a
lieu de dégager les conditions axiomatiques qui assurent les conclusions
tirées dans lesdits n$^{\text{os}}$ .

Nous allons supposer $\underline{P}$ \uline{et} G \uline{plats sur} S, G étant \uline{à fibres}
\var{\inserted[typed]{\var{\xout{de dimen}}{}{typewriter strike-out}{ega25_25_093-v0003}}}{\var{\xout{de dimen}}{}{typewriter strike-out}{ega25_25_093-v0003}}{insertion}{ega25_25_093-v0002}
\uline{géométriquement irréductibles} (pour pourvoir en considérer les points
génériques !), \uline{de dimension} N, \uline{le morphisme} $\underline{P} \hand{\longrightarrow} P$ \uline{est supposé lisse}
\uline{à fibres géométrique\var{\xout{\unsure{s}}}{}{typewriter strike-out}{ega25_25_093-v0004}ment irréductibles de dimension} N-m . Donc le
morphisme $\underline{X} \hand{\longrightarrow} X$ aura les mêmes propriétés. Toutes les propriétés
\unsure{envisaagées|envisagées} \var{\inserted{ici et plus tard,}}{ici et plus tard,}{insertion}{ega25_25_093-v0005} sont stables par \unsure{changment|chqngment} de base sur S, et s'applique\var{\xout{\unsure{r|s}}}{}{typewriter strike-out}{ega25_25_093-v0006}ront
en particulier aux fibres.

Supposons d'abord S=Spec(k). Soit Z une partie fermée de X de
dimension d, alors son image inverse $\underline{Z}$ dans $\underline{X}$ est une partie fermée de
dimension $d + (N-m) = N + d-m$ . Si $d \hand{<} m$, $\underline{Z}$ est de dimension $\hand{<} N$,
donc $\underline{Z} \hand{\longrightarrow} G$ ne peut être dominant, donc si \hand{\unsure{\eta|y}} est le point générique de
G, on a $\underline{Z}_{\hand{\unsure{\eta|y}}} = \hand{\emptyset}$ ; en fait, ce raisonnement montre même (en remplaçant
Z par $\overset{\hand{\rule{1.8em}{0.35pt}}}{f(Z)}$ ) que si dim f(Z) \hand{$<$} m, alors $\underline{Z}_{\hand{\unsure{\eta|y}}} = \hand{\emptyset}$ . On veut une condi-%
tion \var{\inserted[typed]{sur (D)}}{sur (D)}{insertion}{ega25_25_093-v0007} assurant que si dim f(Z) \var{\over{\unsure{$<$}}{$\geq$}}{$\geq$}{writing over}{ega25_25_093-v0008} m, alors $\underline{Z}_{\hand{\unsure{\eta|y}}} \hand{\neq} \hand{\emptyset}$ ; \var{\deleted{et de façon plus
précise que dim $\underline{Z}$ = dim Z - m . Pour la propriété $\underline{Z} \hand{\neq} \hand{\emptyset}$ ,}}{}{deletion}{ega25_25_093-v0009} il

\mspage{93}
